\documentclass[11pt,a4paper]{article}
\usepackage[margin=1in]{geometry}
\usepackage{booktabs}
\usepackage{microtype}
\usepackage{amsmath}
\usepackage[hidelinks]{hyperref}
\usepackage{array}
\usepackage{url}
\sloppy
\setlength{\parskip}{2pt}

\title{Methods supplement: gates, per-round readings, task ledger and provenance inventory\\for the four-scale closed-choice bullwhip study}
\author{Pepijn Frenken\\Independent researcher\\\texttt{pepijnfrenken@mailbox.org}}
\date{22 October 2026}

\begin{document}
\maketitle

\noindent This supplement is the companion to \texttt{paper.tex} and is cited from it as S1--S5.
It exists because the main text had become close to half accounting; the numbers below are the
same numbers, moved, not relaxed. Nothing here changes any number in the paper. Every table names
its source and says whether the source is archived in this repository or only recorded in the
reconciliation note (\texttt{docs/paper/RECONCILIATION-2026-10-22.md}). Where a cited per-run path
is no longer present in the working tree, the table says so instead of implying the file exists.

\section{S1. Gate definitions and frozen thresholds}

\paragraph{Definitions as implemented.} F0 checks the choice grid on 1{,}000 probe states per
regime (regret and top-mass bounds); F1 checks that the 26 label tokens are one contiguous id range
and that the rendered prompt hash is equal across the string renderer and the chat template; F2
freezes the state sample; F3 freezes the collapse-control statistics from that rung's own
\emph{untrained} readout. For the state-dependence statistic we use, over the frozen state sample,
\begin{equation}
  I \;=\; H\!\left(\overline{\pi}\right) \;-\; \overline{H},
\end{equation}
in nats, where $\overline{\pi}$ is the mean of the per-state next-action distributions
$\pi(\cdot\mid s)$ over the 26 label tokens and $\overline{H}$ is the mean of their entropies.
$H_{\mathrm{lo}}$ and $H_{\mathrm{hi}}$ are the 5th and 95th percentiles of the per-state entropy;
$m^{*}$ is the modal share of the oracle argmin letter and $m_{\mathrm{hi}}$ its Wilson upper
bound; $I_{\mathrm{hi0}}$ is the 95th percentile of $I$ over $B{=}1{,}000$ bootstrap resamples.
A trained arm is \emph{collapsed} iff
\begin{quote}
[median entropy $< H_{\mathrm{lo}}$ AND (modal share $> m_{\mathrm{hi}}$ OR
$I \leq I_{\mathrm{hi0}}$)] OR
[median entropy $> H_{\mathrm{hi}}$ AND $I \leq I_{\mathrm{hi0}}$].
\end{quote}
The gate is read at every round end on 256 states and on all 1{,}024 at final evaluation; rounds
are 1-indexed. The stop-loss is terminal after two consecutive collapsing reads; a resume needs a
written override. G1 (pass rule) and G7 (full-vocabulary mass on the 26 label tokens; arm
\texttt{min\_mass} $\geq 0.9$) are stated in the paper's Section 3 and not duplicated here.

\begin{table}[h]
\centering
\small
\caption{Frozen F3 thresholds, to the precision recorded. The 0.6B column is the largest
$H_{\mathrm{lo}}$ of the four, so the 0.6B gate is the \emph{strictest} low-entropy branch.}
\label{tab:s1}
\begin{tabular}{@{}lrrrrl@{}}
\toprule
\textbf{Rung} & $H_{\mathrm{lo}}$ & $H_{\mathrm{hi}}$ & $I_{\mathrm{hi0}}$ & $m_{\mathrm{hi}}$ & \textbf{Archived source} \\
\midrule
0.6B & 1.393 & 2.022 & 0.0970 & 0.149 & pre-flight log line only \\
1.7B & $9.6527\times10^{-5}$ & 1.6313 & 1.4858 & 0.1487 & \emph{no} gate file recovered \\
4B & 0.038026 & 1.910513 & 0.676551 & 0.1487 & \emph{no} gate file recovered \\
8B & 0.000503 & 1.881214 & 1.033805 & 0.1487 & \texttt{gates\_8b.json} \\
\bottomrule
\end{tabular}
\end{table}

\paragraph{Provenance status of Table~\ref{tab:s1}, stated plainly.}
\begin{itemize}
  \item The 0.6B values exist only as the pre-flight job's log line
  (\texttt{results/wave8\_logs/wave8-preflight\_404131.log}):
  \texttt{H\_lo=1.393 H\_hi=2.022 m\_hi=0.149 I\_hi0=0.0970 (m*=0.127)}.
  The archived \texttt{results/wave8/gates.json} contains \texttt{F0}, \texttt{F0\_prereg\_w8\_g26},
  \texttt{F1}, \texttt{S\_nom} and \texttt{S\_nom\_oracle} blocks and \emph{no} \texttt{F3} block.
  \item The 8B values are archived and machine-readable in \texttt{results/wave8/gates\_8b.json},
  which also carries \texttt{F2}: the 26 label ids 32--57, the equal prompt hashes
  (\texttt{55bdc663\ldots}), and an untrained \texttt{min\_mass} of 0.2271.
  \item The 1.7B and 4B values are quoted in the reconciliation note at the same precision as above,
  but \texttt{gates\_1.7b.json} and \texttt{gates\_4b.json} appear in \emph{no} committed revision
  of this repository (\texttt{git log --all --} on both paths is empty) and were not recovered from
  the cluster archive. Those six numbers therefore stand on a document record, not on a gate file.
  This is the weakest link in the 1.7B and 4B gate story and the paper says so.
  \item What is \emph{not} weakened: the readings those thresholds are compared against are
  archived, and Section S2 recomputes the branch verdicts from them. For the 4B, both the
  low-entropy branch's unreachability (minimum median entropy 0.5765 against
  $H_{\mathrm{lo}}{=}0.038$) and the diffuse branch's non-firing (the single read above
  $H_{\mathrm{hi}}$ carries $I{=}0.8714 > I_{\mathrm{hi0}}{=}0.677$) hold with wide margins, so a
  plausible transcription error in the missing gate file would not change the verdict.
  \item The G2 band tiers are a pre-registered \emph{report} tier, not a gate. Both passing 4B
  seeds land in its headline and strong bands at $1.42\times$--$1.56\times$ the archived 5{,}346
  floor; no decision in this paper consumes the tier.
\end{itemize}

\section{S2. Per-round detector readings}

All readings below are copied from archived artifacts, not retyped from prose: the 4B and 1.7B
reads from \texttt{docs/verify/W8-MANIFEST.json} (\texttt{collapse\_4b},
\texttt{stop\_1.7b\_5e5\_seed0}), the 0.6B wave-8 reads from the same manifest
(\texttt{w06\_wave8\_r1\_collapse\_reads}, 16 entries) and the job logs under
\texttt{results/wave8\_logs}. ``C'' in a sequence means the detector reported collapse at that
round read.

\begin{table}[h]
\centering
\small
\caption{4B, both passing seeds, all eight reads. $H$ = median per-state entropy;
$I$ = state-dependence statistic; share = modal argmax share; mass = G7 \texttt{min\_mass}.}
\label{tab:s2a}
\begin{tabular}{@{}llrrrrl@{}}
\toprule
\textbf{Seed} & \textbf{Read} & $H$ & $I$ & \textbf{share} & \textbf{mass} & \textbf{Collapsed} \\
\midrule
0 & R1 (256 st.) & 2.0863 & 0.8714 & 0.1836 & 1.000000 & no \\
0 & R2 (256 st.) & 1.5563 & 0.7063 & 0.4688 & 1.000000 & no \\
0 & R3 (256 st.) & 1.7783 & 0.3459 & 0.7461 & 0.999994 & no \\
0 & final (1{,}024 st.) & 1.7887 & 0.5287 & 0.6514 & 0.999983 & no \\
1 & R1 (256 st.) & 0.5765 & 1.6662 & 0.1953 & 0.999813 & no \\
1 & R2 (256 st.) & 1.7547 & 0.4540 & 0.4102 & 0.999996 & no \\
1 & R3 (256 st.) & 1.7455 & 0.0739 & 0.8398 & 0.999996 & no \\
1 & final (1{,}024 st.) & 1.7239 & 0.3719 & 0.7939 & 0.999996 & no \\
\bottomrule
\end{tabular}
\end{table}

\noindent A note on Table~\ref{tab:s2a}: the reconciliation note cites
\path{results/wave8_r1-4b_lr5e-5_seed{0,1}/round{1,2,3}/collapse_round{n}.json} as the
source of these rows. Those per-run files are not present in the working tree or in any committed
revision; the values themselves are archived in the manifest above and are reproduced here from
it. The same readings under the same thresholds appear in the note, so the discrepancy is a stale
path citation, not a changed number.


\begin{table}[h]
\centering
\footnotesize
\caption{0.6B wave-8, all six arm tasks, all round reads (16 reads). Thresholds:
$H_{\mathrm{lo}}{=}1.393$, $H_{\mathrm{hi}}{=}2.022$, $m_{\mathrm{hi}}{=}0.149$,
$I_{\mathrm{hi0}}{=}0.0970$. Branch: \texttt{low} = low-entropy branch,
\texttt{diff} = diffuse-and-state-blind branch, \texttt{---} = not collapsed.}
\label{tab:s2b}
\begin{tabular}{@{}llrrrcl@{}}
\toprule
\textbf{Arm} & \textbf{Rd} & $H$ & \textbf{share} & $I$ & \textbf{Branch} & \textbf{Seq.} \\
\midrule
\texttt{lr1e-5} s0 & 1 & 0.401 & 0.980 & 0.0991 & low & C \\
\texttt{lr1e-5} s0 & 2 & 2.224 & 0.758 & 0.0196 & diff (reason unrecorded) & C \\
\texttt{lr1e-5} s0 & 3 & 1.760 & 0.391 & 0.5641 & --- & not C \\
\texttt{lr1e-5} s1 & 1 & 0.938 & 0.812 & 0.2046 & low & C \\
\texttt{lr1e-5} s1 & 2 & 1.083 & 0.492 & 0.1606 & low & C \\
\texttt{lr1e-5} s1 & 3 & 1.584 & 0.414 & 0.5879 & --- & not C \\
\texttt{lr1e-5} s2 & 1 & 0.396 & 0.961 & 0.1310 & low & C \\
\texttt{lr1e-5} s2 & 2 & 0.896 & 0.633 & 0.5023 & low & C \\
\texttt{lr1e-5} s2 & 3 & 0.700 & 0.727 & 0.6029 & low & C \\
\midrule
\texttt{lr5e-5} s0 & 1 & 0.826 & 0.508 & 1.0592 & low & C \\
\texttt{lr5e-5} s0 & 2 & 2.087 & 0.695 & 0.2633 & --- & not C \\
\texttt{lr5e-5} s0 & 3 & 0.926 & 0.742 & 0.4367 & low & C \\
\texttt{lr5e-5} s1 & 1 & 0.057 & 0.805 & 0.4679 & low & C \\
\texttt{lr5e-5} s1 & 2 & 0.879 & 0.434 & 0.9350 & low & C \\
\texttt{lr5e-5} s2 & 1 & 0.000 & 0.996 & 0.0309 & low & C \\
\texttt{lr5e-5} s2 & 2 & 1.059 & 0.246 & 1.4025 & low & C \\
\bottomrule
\end{tabular}
\end{table}

\paragraph{Reading Table~\ref{tab:s2b} correctly.} Every round-1 read in the 0.6B wave-8 block has
median entropy below $H_{\mathrm{lo}}{=}1.393$ and modal share above $m_{\mathrm{hi}}{=}0.149$
($H$ spans 0.000--0.938, share 0.508--0.996), so the low-entropy branch fires on all six tasks at
the first read. Two tasks stop there and two more stop after a second collapsing read; the resulting
marker sequences are \texttt{C,C,-} (\texttt{lr1e-5} s0 and s1, both stopped at round 2),
\texttt{C,C,C} (\texttt{lr1e-5} s2, stopped at round 3 when reads 2 and 3 satisfy the
two-consecutive rule), \texttt{C,C} (\texttt{lr5e-5} s1 and s2, stopped at round 2), and
\texttt{C,-,C} (\texttt{lr5e-5} s0), which never meets two \emph{consecutive} collapses and
therefore completed with \texttt{DONE} only rather than being stopped.

\paragraph{One recording anomaly, disclosed.} The \texttt{lr1e-5} s0 round-2 read is marked
\texttt{collapsed: true} with an \emph{empty} \texttt{reasons} list, unlike the other fifteen
entries. The read's own numbers are consistent with the low-entropy branch ($H{=}2.224$ is above
$H_{\mathrm{hi}}{=}2.022$, so the low branch cannot have fired on entropy alone; share 0.758 with
$I{=}0.0196 \leq 0.0970$ is exactly the diffuse-and-state-blind branch, which the detector records
under a different label). We did not reconstruct the missing reason string, and the paper's claim
about this arm rests on the marker sequence and the round-3 read, neither of which depends on it.

\begin{table}[h]
\centering
\small
\caption{1.7B \texttt{lr5e-5} seed 0, the stopped arm: both reads, with branches.}
\label{tab:s2c}
\begin{tabular}{@{}lrrrrl l@{}}
\toprule
\textbf{Round} & $H$ & $I$ & \textbf{share} & \textbf{mass} & \textbf{Branch} & \textbf{Reason} \\
\midrule
1 & $6.2846\times10^{-5}$ & $2.0803\times10^{-6}$ & 1.0000 & 0.9296 & low & share $>m_{\mathrm{hi}}$ \\
2 & 2.0589 & 0.1086 & 0.4453 & $2.1528\times10^{-5}$ & high & diffuse and state-blind \\
\bottomrule
\end{tabular}
\end{table}

\noindent Table~\ref{tab:s2c} is the cleanest single demonstration in the study that the two
detector branches catch different failure modes on the \emph{same} task: round 1 collapses the
output onto one letter with the mass still on the label set (G7 easily passes at 0.9296), and round
2 does the opposite, spreading probability while leaving the label set (G7 fails by five orders of
magnitude). \texttt{stopped\_after\_round: 2}, reason: ``collapse gate tripped at two consecutive
round reads (WAVE8 \S6.5)''. The paper reports the $H$ and $I$ comparisons against this rung's
$H_{\mathrm{lo}}{=}9.6527\times10^{-5}$ and $H_{\mathrm{hi}}{=}1.6313$; both are in
Table~\ref{tab:s1} and carry the provenance caveat noted there.


\section{S3. The full 24-row task ledger}

One row per planned task (scale $\times$ recipe arm $\times$ learning-rate arm $\times$ seed).
Dispositions use the precedence stated in the paper: a task stopped by the collapse stop-loss is
counted as \emph{stopped} even where a later resume produced a completed read, so the five classes
below are disjoint. ``Matched final eval'' means a paired $n{=}30$ noisy-argmax generation against
that rung's matched untrained control.

\begin{table}[h]
\centering
\footnotesize
\caption{All 24 planned tasks. $n{=}30$ paired paths per evaluated row. Ctl = matched untrained
control median.}
\label{tab:s3}
\begin{tabular}{@{}p{2.35cm}p{2.0cm}p{2.6cm}p{3.5cm}p{3.2cm}@{}}
\toprule
\textbf{Task} & \textbf{Sequence} & \textbf{Disposition} & \textbf{Measured read} & \textbf{Evidence} \\
\midrule
0.6B w7 SFT s0 & --- & evaluated, failing & 20{,}594 vs ctl 11{,}569; 0/30 wins & per-run JSON \\
0.6B w7 SFT s1 & --- & evaluated, failing & 25{,}721; 0/30 & per-run JSON \\
0.6B w7 SFT s2 & --- & evaluated, failing & 25{,}508; 0/30 & per-run JSON \\
0.6B w7 xgrad s0 & --- & evaluated, failing & 37{,}384; 0/30 & per-run JSON \\
0.6B w7 xgrad s1 & --- & evaluated, failing & 32{,}516; 0/30 & per-run JSON \\
0.6B w7 xgrad s2 & --- & evaluated, failing & 27{,}652; 0/30 & per-run JSON \\
\midrule
0.6B w8 lr1e-5 s0 & C,C,- & stopped (resumed; 2 markers) & R3 $H{=}1.760$ clean; final G7 $3.2\times10^{-4}$ & log + marker \\
0.6B w8 lr1e-5 s1 & C,C,- & stopped (resumed, override) & no final eval & log + marker \\
0.6B w8 lr1e-5 s2 & C,C,C & stopped at R3 & no final eval & log + marker \\
0.6B w8 lr5e-5 s0 & C,-,C & trained, no final eval & \texttt{DONE} only; band 16{,}035 & log + band record \\
0.6B w8 lr5e-5 s1 & C,C & stopped at R2 & no final eval & log + marker \\
0.6B w8 lr5e-5 s2 & C,C & stopped at R2 & no final eval & log + marker \\
\midrule
1.7B lr1e-5 s0 & not C & evaluated, failing & 21{,}468 vs ctl 11{,}508; 0/30 & per-run JSON \\
1.7B lr1e-5 s1 & one trip & evaluated, failing & 43{,}660; 0/30; G7 fail & per-run JSON \\
1.7B lr1e-5 s2 & not C & evaluated, failing & 71{,}270; 0/30 & per-run JSON \\
1.7B lr5e-5 s0 & C,C & stopped at R2 & both branches (Table~\ref{tab:s2c}) & manifest \\
1.7B lr5e-5 s1 & --- & unattempted & never submitted & submit ledger \\
1.7B lr5e-5 s2 & --- & unattempted & never submitted & submit ledger \\
\midrule
4B lr5e-5 s0 & not C & evaluated, passing & 8{,}315 vs ctl 21{,}578; 30/30 & verdict doc; raw
JSONL lost (S5) \\
4B lr5e-5 s1 & not C & evaluated, passing & 7{,}612; 30/30 & verdict doc; raw JSONL lost (S5) \\
\midrule
8B lr5e-5 s0 & n/a & Oct trained, no eval; Nov band read only & Oct band 35{,}289; Nov band 93{,}603$\to$10{,}097, 30/30 & verdict doc; Nov \texttt{.costs.jsonl} \\
8B lr5e-5 s1 & n/a & Oct trained, no eval; Nov band read only & Oct band 12{,}444; Nov band 93{,}603$\to$15{,}355, 28/30 & verdict doc; Nov \texttt{.costs.jsonl} \\
8B lr1e-5 s0 & n/a & Oct trained, no eval; Nov attempt lost & Oct band 71{,}855; Nov leg died in round 3 & verdict doc \\
8B lr1e-5 s1 & n/a & Oct trained; band-only & band 22{,}083 (pooled 46{,}969), $0.21\times$ control & verdict doc; adapter on Modal volume \\
\bottomrule
\end{tabular}
\end{table}

\paragraph{Totals, reconciled.} Planned 24. Attempted with a trained checkpoint 22 (the two
unattempted are 1.7B \texttt{lr5e-5} seeds 1--2). With a matched paired final evaluation 11 (six
wave-7, three 1.7B 1e-5, two 4B). Without one 13 (five stopped wave-8 0.6B, one \texttt{DONE}-only
wave-8 0.6B, one stopped 1.7B 5e-5, two unattempted, four 8B). Six distinct tasks carry a persisted
\texttt{STOPPED\_collapse} marker, of which two were later resumed. The five disposition classes sum
to 2 + 6 + 5 + 9 + 2 = 24 exactly as stated in the paper's Table~1, and the six-marker count is a
\emph{separate flag} over the same rows, not a sixth class: this is the overlap the paper discloses
rather than denies.

\section{S4. The registered denominator text, and both estimands}

\paragraph{Verbatim.} \texttt{docs/waves/WAVE8\_SPEC.md} \S6.3, committed 2026-10-07 (commit
\texttt{7d3260b}), before any wave-8 run:

\begin{quote}
``A seed passes iff all three hold; the arm needs $\geq$ 2/3 seeds.''
\end{quote}

\noindent The sentence does not define how a planned seed that never produces an evaluated
checkpoint enters the fraction. The registered treatment is therefore \emph{underspecified}; it is
not amended retrospectively, and the paper prints two readings instead.

\begin{table}[h]
\centering
\small
\caption{Both denominator readings, at the rule's actual unit (the learning-rate arm).}
\label{tab:s4}
\begin{tabular}{@{}p{2.4cm}p{2.3cm}p{3.4cm}p{4.0cm}@{}}
\toprule
\textbf{Rung} & \textbf{Arm} & \textbf{Evaluated reading} & \textbf{Planned reading} \\
\midrule
4B & 5e-5 & $2/2$ pass & $2/2$ pass \\
1.7B & 1e-5 (executed) & $0/3$ & $0/3$ \\
1.7B & 5e-5 (planned) & undefined (0 evaluated) & $0/3$ \\
1.7B & rung level & undefined & $0/6$ \\
0.6B w8 & 1e-5 & undefined (0 evaluated) & $0/3$ \\
0.6B w8 & 5e-5 & undefined (0 evaluated) & $0/3$ \\
0.6B w8 & pooled & undefined & $0/6$ \\
8B & both & not computable (no matched control) & not computable \\
\bottomrule
\end{tabular}
\end{table}

\noindent Two consequences the paper states and this table makes arithmetic: the only positive rung
is unchanged under both readings, and the arms with zero evaluated seeds are \emph{undefined}
rather than \emph{failing}. A third consequence belongs to the selection step and is easy to miss:
the rule ranks arms by pooled mean validation cost \emph{over seeds}, so a missing seed corrupts arm
selection independently of the pass fraction. At 8B no seed is missing and selection was computable
(pooled 23{,}867 vs 46{,}969, \texttt{winner:} \texttt{lr5e-5}); at 1.7B the selection record holds
only the executed arm (pooled 48{,}311 over three seeds, \texttt{winner:} \texttt{lr1e-5}) because
the planned 5e-5 arm produced one trained-but-stopped seed with no validation read; at wave-8 0.6B
the band record holds one seed per arm (90{,}448 and 16{,}035, \texttt{winner:}
\texttt{lr5e-5}), so no pooled-mean selection was performed there either.


\section{S5. Provenance inventory}

\subsection*{Commit anchors}

\begin{table}[h]
\centering
\small
\caption{Anchors cited by the paper, with what each one pins. Times UTC unless stated.}
\label{tab:s5a}
\begin{tabular}{@{}p{1.7cm}p{3.6cm}p{7.0cm}@{}}
\toprule
\textbf{Commit} & \textbf{Date / time} & \textbf{What it pins} \\
\midrule
\texttt{7d3260b} & 2026-10-07 13:01:04Z & \texttt{docs/waves/WAVE8\_SPEC.md} with F0--F3, G1, G2, G7, the
seed-pooling rule and the failing-grid record. Byte-identical spec text at the two later commits below \\
\texttt{4b8bffb} & 2026-10-07 14:55:36Z & cluster kit, byte-equal to the spec text at \texttt{7d3260b} \\
\texttt{b61cea6} & 2026-10-07 14:57:38Z & spec status header + appended changes log only; training arrays checked this out \\
\texttt{d720701} & 2026-10-08 08:05:31Z & later checkout of the same spec text by the training arrays \\
\midrule
& 2026-10-08 22:24:06Z & first wave-8 GPU job \texttt{405143} starts, 7.5 h after the kit commit \\
\bottomrule
\end{tabular}
\end{table}

\noindent The provenance claim in the paper is narrow and verifiable: \texttt{sha256sum} of
\texttt{git show <sha>:docs/waves/WAVE8\_SPEC.md} is identical at \texttt{7d3260b}, \texttt{b61cea6}
and \texttt{d720701}; \texttt{git merge-base --is-ancestor} confirms \texttt{7d3260b} is an ancestor
of both training HEADs; and \texttt{git diff 7d3260b b61cea6 -- docs/waves/WAVE8\_SPEC.md} shows only
the status header and an appended changes log. No rule text changed after the freeze.

\subsection*{Job and lane inventory}

\begin{table}[h]
\centering
\small
\caption{Compute lanes and job identifiers.}
\label{tab:s5b}
\begin{tabular}{@{}p{2.0cm}p{2.6cm}p{8.0cm}@{}}
\toprule
\textbf{Rung} & \textbf{Lane} & \textbf{Identifiers} \\
\midrule
0.6B -- 4B & TU/e Umbrella, L4 (24 GB) & pre-flight \texttt{404131}; wave-8 training arrays
\texttt{405143}, \texttt{405887}, \texttt{406726}, \texttt{407006} (per-task array elements) \\
8B, seed 0 & Snellius, A100 & \texttt{28414248}, \texttt{28418615} \\
8B, seed 1 & Modal, A100 & two legs after the budget amendment; job identifiers in the 8B verdict doc \\
\bottomrule
\end{tabular}
\end{table}

\noindent Array-element start times from the manifest's \texttt{job\_timeline}: \texttt{405143\_0}
2026-10-08T00:24:02 local, \texttt{405887\_0..5} from 09:35:08, \texttt{406726\_0..2} from 12:11:35,
\texttt{407006\_0..2} from 13:19:09. The two later arrays are the resumed runs whose markers are
discussed in S2; the first is the one that produced most of the round-1 collapse trips.

\subsection*{Code and configuration}

Implemented modules named by the paper: \texttt{agent\_bullwhip/engine.py} (four-tier beer-game
simulator), \texttt{decision\_gates.py} (F0--F3, collapse detector, permutation check),
\texttt{make\_w8\_g26b()} (the 26-value frozen grid), \texttt{verify.py} (gate orchestration and
pristine-block gate), \texttt{gates.py}. Per-rung frozen gate files: see Table~\ref{tab:s1}. The
per-provider library stack is recorded in the cluster environment specification and was
\textbf{not} runtime-verified in preparing this paper, so it is named as a declaration, not a
measurement.

\subsection*{Evidence availability, per rung}

\begin{table}[h]
\centering
\small
\caption{What an independent reader can actually re-read today, and what they cannot.}
\label{tab:s5c}
\begin{tabular}{@{}p{1.5cm}p{4.4cm}p{6.4cm}@{}}
\toprule
\textbf{Rung} & \textbf{Available} & \textbf{Not available} \\
\midrule
0.6B w7 & per-run JSON for all six arm--seed rows; matched untrained control & --- \\
0.6B w8 & job logs with every round read and marker; \texttt{w8-band-record.json};
sixteen reads in the manifest & arm directories and the two curriculum-band files (removed by the
repository's results slim-down, not by this paper) \\
1.7B & \texttt{gates\_1.7b.json}, re-frozen 2026-11-20 and matching the registered values; stop
record for the 5e-5 arm; November rebuild per-run JSONL for the three 1e-5 seeds and the untrained
controls (S6) & October per-run JSONL for the three evaluated seeds, either generation; any
final-eval files for the stopped arm (they never existed) \\
4B & \texttt{gates\_4b.json}, re-frozen 2026-11-20; November rebuild per-run JSONL for both seeds,
permutation JSONs and both generations (S6) & October raw paired JSONLs for both seeds; October
permutation JSONs \\
8B & verdict doc with leg-by-leg timings and band means; the lr1e-5 s1 checkpoints (adapter
30.7 MB) and band table on the Modal volume; November readouts for both generations, the untrained-8B
control and its band (S7) & selection JSON; the other three arms' adapters; October raw JSONL of
either generation \\
\bottomrule
\end{tabular}
\end{table}

\paragraph{Closing statement on scope.} This supplement exists to let a reader check the paper's
numbers rather than trust its prose. Four gaps were open when the October results were written: the
six 1.7B and 4B threshold values with no gate file recovered, the 0.6B wave-8 per-run directories
removed by the repository's own slim-down, the 8B raw evaluations never recovered from either lane,
and one missing detector reason string in the 0.6B wave-8 block. Their status now differs, and we
state it per item rather than as a single caveat. The 1.7B and 4B thresholds are no longer
unchecked: the November rebuild re-froze both from the pinned checkpoints and every value matches
the registered figures within $3.3\times10^{-4}$ relative (S6, Table~\ref{tab:s6b}), so the gate
values are confirmed by reproduction even though the October files were not recovered. The 0.6B
per-run directories, the 8B raw evaluations and the reason string remain missing. The October
numbers themselves are \emph{not} recomputable from raw data: the 1.7B and 4B per-run JSONLs are
gone, which is why S6 reports an independent rerun rather than a reconstruction. That rerun is what
now backs the paper's two verdicts.

\section{S6. Independent rebuild, November 2026}\label{sec:s6}

The wave-8 per-run artifacts were lost with the cluster scratch archive (S5). This section reports an
independent rerun of the same protocol from the committed kit on the same pinned checkpoints, executed
on 2026-11-19 and 2026-11-20 on the TU/e Umbrella L4 lane as jobs \texttt{766006}, \texttt{766026},
\texttt{766027}, \texttt{766028}, \texttt{766029}, \texttt{766030}, \texttt{766031}, \texttt{775521}
and \texttt{775522}. All completed with exit 0 and no run was stopped: one round-1 read tripped the
collapse detector (1.7B seed 1) and is recorded below, and training continued past it by design at
that rung. The method is unchanged: same gate freeze, same training recipe, same evaluation
generations, same pins.

\subsection*{Why the rebuild is admissible as a reproduction}

It writes its own gate files before training and refuses to start when a statistic drifts. Table
\ref{tab:s6b} is that check, byte-read from \texttt{results/wave8/gate\_verification.json}, whose
\texttt{ok} field is \texttt{true} across all eight checks. Because the gate values were fixed in
\emph{October}, the match is a genuine prediction rather than a re-fit.

Second, the readout side is deterministic to the byte. Two independent evaluation jobs of the same 4B
checkpoints produced identical files: \texttt{sha256} prefixes \texttt{24821f8c4eb2}
(\texttt{w8-r1-4b-seed0-argmax-noisy.jsonl}), \texttt{90663e3f079e}
(\texttt{w8-r1-4b-seed1-argmax-noisy.jsonl}) and \texttt{77b8590bec15}
(\texttt{w8-r1-4b-seed0-sampled-step.jsonl}). Any difference from October therefore sits in training,
not in the evaluation.

\begin{table}[h]
\centering
\small
\caption{Re-frozen gate statistics against the values registered before the October runs. Tolerance
$10^{-3}$ relative. Read from \texttt{results/wave8/\allowbreak gate\_verification.json}.}
\label{tab:s6b}
\begin{tabular}{@{}llrrr@{}}
\toprule
\textbf{Rung} & \textbf{Field} & \textbf{Registered} & \textbf{Re-frozen} & \textbf{Rel.\ diff} \\
\midrule
1.7B & $H_{\mathrm{lo}}$ & $9.6527\times10^{-5}$ & $9.6527\times10^{-5}$ & $4.6\times10^{-6}$ \\
1.7B & $H_{\mathrm{hi}}$ & 1.6313 & 1.6313 & $1.9\times10^{-6}$ \\
1.7B & $I_{\mathrm{hi0}}$ & 1.4858 & 1.4858 & $1.8\times10^{-5}$ \\
1.7B & $m_{\mathrm{hi}}$ & 0.1487 & 0.1487 & $3.3\times10^{-4}$ \\
4B & $H_{\mathrm{lo}}$ & 0.038026 & 0.038026 & $9.8\times10^{-6}$ \\
4B & $H_{\mathrm{hi}}$ & 1.910513 & 1.910513 & $2.5\times10^{-7}$ \\
4B & $I_{\mathrm{hi0}}$ & 0.676551 & 0.676551 & $7.1\times10^{-7}$ \\
4B & $m_{\mathrm{hi}}$ & 0.1487 & 0.1487 & $3.3\times10^{-4}$ \\
\bottomrule
\end{tabular}
\end{table}

\subsection*{Untrained controls}

The noisy-argmax control medians are 21{,}578 (4B) and 11{,}508 (1.7B), identical to the values in
the main text, and the 1.7B noisy mean is 11{,}405, also identical. The curriculum-band control means
are 116{,}948.33 for 4B, which reproduced to the byte on a second job 21 hours later, and 20{,}154.80
for 1.7B. Files: \texttt{results/eval/w8-untrained-\{1.7b,4b\}-argmax-noisy.jsonl} and
\texttt{results/wave8\_val/untrained-curriculum-\{1.7b,4b\}.jsonl}.

\subsection*{G1 on the rebuild}

G1 is evaluated on the generation the spec names for it, paired noisy-argmax $n{=}30$
(\texttt{WAVE8\_SPEC.md} \S G1): median cost reduction $\geq 15\%$ in both variants, $\geq 20/30$
paired wins, one-sided sign test $p<0.02$. Demand paths depend only on the seed base and are
model-independent, so control and trained rows pair by run index. No pair tied.

\begin{table}[h]
\centering
\small
\caption{G1 on the rebuild, paired noisy-argmax $n{=}30$, against each rung's own untrained control.
Odds ratio is the ratio of medians; the per-path ratio is the median of the 30 individual ratios.}
\label{tab:s6a}
\begin{tabular}{@{}lrrrrrll@{}}
\toprule
\textbf{Arm} & \textbf{Wins} & \textbf{Med.\ ctl} & \textbf{Med.\ trained} & \textbf{Ratio} & \textbf{Per-path} & \textbf{$p$} & \textbf{G1} \\
\midrule
4B seed 0 & 30/30 & 21{,}578 & 9{,}346 & 0.43 & 0.40 & $9.3\times10^{-10}$ & pass \\
4B seed 1 & 25/30 & 21{,}578 & 17{,}958 & 0.83 & 0.75 & $1.6\times10^{-4}$ & pass \\
\midrule
1.7B seed 0 & 7/30 & 11{,}508 & 27{,}706 & 2.41 & 2.28 & --- & fail \\
1.7B seed 1 & 0/30 & 11{,}508 & 79{,}094 & 6.87 & 7.88 & --- & fail \\
1.7B seed 2 & 0/30 & 11{,}508 & 94{,}887 & 8.25 & 8.23 & --- & fail \\
\bottomrule
\end{tabular}
\end{table}

Both verdicts reproduce: 4B clears G1 on both seeds (2/2, so the arm passes), and 1.7B fails on all
three (0/3, adverse). The magnitudes are not identical. The 4B seed-1 median is 17{,}958 rather than
7{,}612 and its win count 25/30 rather than 30/30; it clears the 15\% reduction bar by two points on
the ratio of medians. The 1.7B regressions are larger than published (2.41, 6.87, 8.25 against 1.87,
3.79, 6.19), and its seed 0 gets 7/30 wins where the paper records 0/30.

\subsection*{What differs, and where}

Training outcomes are not stable across the two executions even though the readout is. On the
curriculum band, 4B moves from 10{,}547.8 and 14{,}681.8 (October) to 9{,}355.1 and 19{,}385.3
(November), pooled 12{,}614.8 to 14{,}370.2. The 1.7B band moves from 47{,}125 / 36{,}734 / 61{,}073
to 28{,}825 / 86{,}546 / 86{,}595, and the seed ordering inverts: October's cheapest 1.7B seed is
seed 1, November's is seed 0. Since evaluation is byte-reproducible, this variance belongs to training,
and it is reported rather than smoothed: the paper's per-seed numbers are single-execution facts, and
the seeds are more dispersed than the October pair suggests.

\subsection*{Mechanism, unchanged}

Median entropy per training round, read from
\texttt{results/wave8\_r1-*/round*/collapse\_round*.json}: 4B holds 1.885 / 1.661 / 1.486 (seed 0) and
1.525 / 1.493 / 1.461 (seed 1). 1.7B falls to $2.16\times10^{-2}$ / $2.89\times10^{-3}$ /
$3.20\times10^{-6}$ (seed 0), $1.02\times10^{-1}$ / $3.00\times10^{-2}$ / $2.50\times10^{-4}$ (seed 2),
and seed 1 trips the collapse gate in round 1 with modal argmax share 1.0 before recovering to 0.243
and falling back to $3.2\times10^{-3}$.

One dissociation belongs here because a simpler reading would be wrong: 1.7B seed 0 ends at the lowest
median entropy of any arm and is still the cheapest 1.7B arm (median 27{,}706 against 79{,}094 and
94{,}887). Degeneracy and cost are not the same axis. A near-deterministic policy is cheap when its
single action happens to be the canonical one. The defensible statement is that the 1.7B loses the
policy and its cost rises; ``costs more because it collapses'' is too strong for these data.

\subsection*{Two caveats carried forward}

The untrained 1.7B read has \texttt{min\_mass} 0.60945 on the band and 0.85869 (argmax-step), 0.543
(sampled-step) and 0.82305 (argmax-noisy) on the single runs, all below the C0 threshold of 0.99. The
1.7B comparison therefore rests on a control whose readout partly leaves the 26-letter vocabulary. The
direction is unaffected; the mass caveat belongs in the limitations.

Second, one defect was found and fixed in the rebuild's own kit rather than in the protocol: the
curriculum-band control tag carried no rung suffix, so the 1.7B and 4B evaluation jobs wrote the same
path and the last writer won. Both files are now written with the rung in the name, and both are
reported above. It is stated because a duplication of that kind would have silently deleted half the
evidence for a second time.


\section{S7. The 8B rung, November 2026}\label{sec:s7}

The body reports 8B as evaluable on one arm, the October survivor whose paired generation exists,
and reports the selected learning rate as band-level favourable but unestablished under G1. This
section carries the artifacts behind both sentences in the order they were produced: the one adapter
that survived October, the two November re-training legs at $5\times10^{-5}$, and the losses,
including a third November leg that died with its sandbox, the eval traces a transfer corrupted, and
a mislabelled first read of the November numbers that this section now corrects. The sequence is the
evidence, so the failures are kept.

\subsection*{The November re-training legs at the selected learning rate}

Two legs, lr $5\times10^{-5}$, seeds 0 and 1, run on 2026-11-22 on one RTX PRO 6000 Blackwell
card with 95.6\,GiB usable, 20 CPUs, torch 2.11.0+cu130, transformers 4.55.4 and peft 0.21.1.
Configuration identical for both: \texttt{decision\_rounds run --arm r1}, three rounds of 170
steps, batch 32, maximum length 1024, gradient checkpointing, save every 10 steps, eight rollout
processes, gates \texttt{results/wave8/gates\_8b.json}. Frozen inputs were checked before launch:
\texttt{S\_nom.jsonl} at \texttt{f651fc58\ldots6aaa8} and the gate file at \texttt{b6194407\ldots1f90}.

\begin{center}\small
\begin{tabular}{lrrrrrr}
\hline
seed & wall & band mean (min--max) & min mass & wins & ratio & $p$ \\
\hline
0 & 62\,min & 10{,}892 (7{,}168--19{,}524) & 0.99997 & 30/30 & 0.108 & $9.3\times10^{-10}$ \\
1 & 62\,min & 15{,}411 (13{,}059--19{,}586) & 0.99973 & 28/30 & 0.164 & $4.3\times10^{-7}$ \\
\hline
\end{tabular}
\end{center}

Both read against the \emph{curriculum} control they were run against,
\texttt{results/wave8\_val/untrained-curriculum-8b.jsonl}, median 93{,}602.5 and mean 104{,}676.07;
wins and the sign test are per paired run index. Reproduce with
\texttt{scripts/wave8\_g1\_read.py results/wave8\_val/untrained-curriculum-8b.jsonl
results/wave8\_val/w8-val-r1-8b-lr5e-5-seed\{0,1\}.costs.jsonl}. The two seeds sit
$1.52\times$ apart on the same control.

\textbf{This is not the G1 estimand, and a first version of this section got that wrong.} G1 is
frozen on the paired \emph{noisy-argmax} generation. The November legs were run with
\texttt{--pattern curriculum}, so their 30 runs share demand paths with the curriculum control and
none with the noisy control (\texttt{incoming\_now} vectors: 12 of 12 parseable seed-0 rows match
curriculum, 0 of 12 match noisy). An earlier draft of this section paired them against
\texttt{results/eval/w8-untrained-8b-argmax-noisy.jsonl} and reported 29/30 wins with
$p=2.9\times10^{-8}$; that pairing is void because the two sides never share a demand path. The
prompt SHA is the same in all four files precisely because the rendered prompt is hashed
independently of the demand draw, so the hash cannot certify a matched pair. The selected rate's G1
outcome is therefore unestablished; what exists is a favourable band comparison, which is what the
body now reports.

\textbf{Collapse gates.} No round tripped the stop-loss on either seed; the $I$ statistic reads
0.5017, 0.5103 and 1.0998 across seed 0's three rounds. The gate JSON files record only $I$ and
write \texttt{medH} and \texttt{modal\_share} as null, while the console printed 0.001 and 0.895 for
round 1 and 0.011 and 0.699 for round 3. A reader opening the files sees the nulls, so the printed
values are quoted here rather than inferred.

\textbf{What was lost, in this attempt.} A third leg, $10^{-5}$ seed 0, started after seed 1 and
was in its third round when the sandbox was reclaimed; it wrote no validation read and is not
evaluable. The LoRA adapters of both completed legs went down with the instance. The full-trace
evaluation JSONLs were pulled as a 750\,KB base64 payload inside a cell print and landed corrupt,
12 of 30 rows parseable and the file ending mid-object; they are removed rather than kept, and the
30-row cost extracts that G1 actually consumed are committed in their place. The defect is named
because it will recur: printed blobs above a few kilobytes do not survive that channel intact.

\subsection*{What survived and where}

Three of the four 8B legs lost their adapters with the archives they ran in. The 5e-5 seed-1 leg's
archive on the run volume is \texttt{leg\_5e-5\_seed1.zip} at \textbf{22 bytes}: the archival step
crashed before writing. The $10^{-5}$ seed-1 leg is the exception. Its directory holds
\texttt{round\_state.json}, a \texttt{DONE} marker, \texttt{collapse\_reads.jsonl}, three rounds of
checkpoints, and a validation-band table
(\texttt{w8-val-r1-8b-lr1e-5-seed1.summary.json}: 30 runs, mean 22{,}083.37, \texttt{min\_mass}
0.99205). The adapter at \texttt{checkpoint-169/adapter\_model.safetensors} is 30{,}709{,}192 bytes,
and its \texttt{state.json} records \texttt{lr} 1e-5, \texttt{seed} 1, \texttt{step} 169, model
\texttt{Qwen/Qwen3-8B}, saved 2026-10-22T14:06:37Z. That one file is the difference between an
unevaluable rung and a partial one.

The gate file \texttt{results/wave8/gates\_8b.json} (1{,}497 bytes, frozen 2026-10-21,
\texttt{sha256} \texttt{b6194407e03a}) carries the diagnostic thresholds: \texttt{H\_lo}
0.0005025333805207596, \texttt{m\_hi} 0.14874833921755284, \texttt{min\_mass} 0.22705760598182678 on
the frozen diagnostic state sample.

\subsection*{Two failed attempts, and why they are recorded}

\textbf{Job \texttt{780799} failed in five seconds} with \texttt{FATAL: results/wave8/gates\_8b.json
missing}. The file existed on one disk and had never been force-added to git, so the cluster checkout
never held the 8B gate freeze at all. It is now committed (\texttt{3f71c90}) and both copies hash
\texttt{b6194407e03a}. A paper whose 8B row depends on that file was one disk failure away from losing
its thresholds.

\textbf{Job \texttt{780926} hit the memory wall.} Seven of nine steps raised
\texttt{torch.cuda.OutOfMemoryError} at 21.43\,GiB of a 22.06\,GiB L4: the runner batches every demand
path in lockstep, so 30 paths $\times$ 4 roles is 120 prompts in one forward. Two steps landed first
and are reported below because they are valid.

\textbf{The fix and its scope.} \texttt{W8\_READOUT\_CHUNK} slices the prompt list before the forward
pass. The readout is per prompt (left-padded, one forward per group, logits read at the last real
position, mass computed per prompt), so chunking changes the grouping and nothing else. Four prompts
at a time is the footprint the single-path runs had already proved on this card. It is an
infrastructure parameter, not a change to a gate, a threshold or the pooling rule; the 0.6B, 1.7B and
4B rungs were not re-run with it and are unaffected.

\subsection*{Readouts, job \texttt{781221}, all steps exit 0, 1h20m43s}

\begin{table}[h]
\centering
\small
\begin{tabular}{@{}lrrrr@{}}
\toprule
\textbf{arm} & \textbf{generation} & \textbf{runs} & \textbf{mean cost} & \textbf{min mass} \\
\midrule
untrained & argmax-step & 1 & 93{,}060.00 & 0.84032 \\
untrained & sampled-step & 30 & 109{,}183.47 & 0.79728 \\
untrained & argmax-noisy & 30 & 106{,}448.03 & 0.78885 \\
adapted & argmax-step & 1 & 8{,}376.00 & 0.99986 \\
adapted & sampled-step & 30 & 9{,}385.13 & 0.99706 \\
adapted & argmax-noisy & 30 & 14{,}525.43 & 0.99810 \\
\bottomrule
\end{tabular}
\caption{8B readouts from \texttt{results/eval/w8-*-8b-*.summary.json}. The adapted arm is lr $10^{-5}$
seed 1.}
\label{tab:s7a}
\end{table}

\subsection*{The frozen G1 read}

\begin{table}[h]
\centering
\small
\begin{tabular}{@{}ll@{}}
\toprule
\textbf{quantity} & \textbf{value} \\
\midrule
paired paths & 30 (no ties) \\
control median, noisy argmax & 89{,}633.00 \\
arm median & 14{,}594.50 \\
median ratio & 0.1628 (83.72\% reduction) \\
median of per-path reductions & 0.1569 (84.31\%) \\
paired wins & 29/30, one against \\
one-sided sign test & $p = 2.89\times10^{-8}$ \\
\bottomrule
\end{tabular}
\caption{G1 at 8B, lr $10^{-5}$ seed 1 against the matched untrained-8B control. Reproduce with
\texttt{python3 scripts/wave8\_g1\_read.py results/eval/w8-untrained-8b-argmax-noisy.jsonl
results/eval/w8-r1-8b-seed1-argmax-noisy.jsonl}.}
\label{tab:s7b}
\end{table}

Every gate passes: the 15\% bar on both statistics, 20/30 wins, $p<0.02$, and the 22/30 the abstract
text asks for. On the curriculum band the same arm reads 22{,}083.37 against an untrained-8B band of
104{,}676.07, a 0.211 ratio, so band and generation agree in direction.

The selected 5e-5 arm has no surviving checkpoints and therefore no paired generation, but it does
have document-recorded band means, and against the same untrained band those sit at $0.337$
(35{,}289) and $0.119$ (12{,}444), pooled $0.228$. That gives the selected arm a band-level direction
without a retrain. It is not a substitute for G1: the provenance is the October verdict document
rather than per-run files, and the control's mass caveat applies to it as much as to the other arm.

\subsection*{Caveats that belong with the number}

\textbf{One arm, and not the selected one.} The frozen rule selects 5e-5 by pooled band mean
(23{,}867 against 46{,}969), and that arm's adapters are gone. The passing arm is the one the rule did
not select. No post-hoc arm switch is taken, so the rung is one arm passing and the selected arm
unevaluated, which is weaker than an 8B verdict.

\textbf{The control's readout leaves the vocabulary.} The untrained 8B resolves 0.63050 of its mass on
the 26 letters on the band and 0.84032, 0.79728 and 0.78885 on the three generations, all below the
C0 threshold of 0.99 and far below the adapted arm's 0.99706 to 0.99986. The frozen gate file's
diagnostic figure is 0.22706, a different state sample, which is why that artifact could not serve as
the control. The comparison therefore rests on a control that is partly outside the closed choice set.
The direction is not in doubt at 29/30 with an 83.7\% reduction; the magnitude carries that caveat.

\textbf{Permutation invariance, in context.} The adapted arm's final read is 0.0742, 38 of 512
cases invariant (\texttt{results/eval/w8-r1-8b-seed1-perm.json}), against a registered gate of 0.9.
No arm passes that gate on any rung: 4B seed 0 is 0.0020 and seed 1 0.0273, 1.7B is 0.1172, 0.1543 and
0.1562, and all wave-7 arms are 0.0000. The 8B arm therefore sits fourth among the trained arms, and
an earlier draft's 0.0781 with 40 of 512 came from the failed first job rather than the final one.
\texttt{WAVE8\_SPEC.md} makes invariance below 0.05 a stop only in combination with a tripped collapse
gate. It is reported so that a reader meeting the raw file does
not mistake a universal property for an 8B failure.

The 4B figure in this table comes from the November generation and does not agree with the October
one: \texttt{W8-MANIFEST.json} records 24/512 (0.046875) and 20/512 (0.0390625) for the same two
seeds, because the November eval wrote files of the same name and the manifest kept the earlier
read. Both cohorts pass the verdict that no arm reaches 0.90, and the 24-fold gap is itself the
usable result: the permutation read is not stable across generations of one rung, so it supports a
withheld interface claim and nothing finer.

\textbf{Hardware.} The 8B readouts ran on the same L4 lane as the 0.6B--4B rungs, not on the A100 lanes
the 8B \emph{training} used. Same code, same pinned checkpoint, same demand paths.

\end{document}
